\chapter{Free Particles, Momentum Space, and the Classical Limit}

These notes follow Leonard Susskind's lecture closely, keeping the order of motivation, the pauses for interpretation, and the later return to classical mechanics. Transcript and frame curation by Leonard Susskind, LazyingArt LLC, and Video2Book are used here only to preserve the visible board work and the sequence of the spoken argument. The lecture begins far from the Schr\"odinger equation, with the historical pressure that made quantum mechanics unavoidable, and only then returns to the mathematics of a particle on a line.

\section{From Quantum Crisis to Quantum First Principles}

The lecture opens by clearing away the temptation to think of quantum mechanics as a small correction to classical reasoning. Michelson--Morley is part of that background, but the deeper crisis for the present lecture is blackbody radiation. The old wave theory of light said that electromagnetic radiation could occur at arbitrary wavelength, and each wavelength behaved mathematically like a harmonic oscillator. Statistical mechanics then said that in thermal equilibrium each oscillator should carry the same average energy,
\begin{equation}
E_{\text{mode}} = kT.
\end{equation}
But a cavity contains infinitely many modes. If every allowed mode carries the same nonzero energy, then the total energy diverges. This is the ultraviolet catastrophe.

Planck's move was to write
\begin{equation}
E = n h \nu,
\end{equation}
but, as the lecture emphasizes, he did not yet interpret this correctly as a statement about the radiation in the cavity itself. In Planck's original picture, the quantization belonged to the wall oscillators that emitted the radiation. Einstein's deeper step was to say that the radiation itself comes in quanta. Then the high-frequency modes are suppressed: if
\begin{equation}
h\nu > kT,
\end{equation}
the classical equipartition estimate would assign less than one quantum to that mode, and quantum mechanics does not allow that. Less than one quantum means no quanta.

From there the lecture moves through the slow acceptance of photons, the role of Compton scattering, and the emergence of genuine probabilistic thinking. The point is not antiquarian. Susskind wants the logical direction to be clear before the mathematics resumes: quantum mechanics is the more fundamental theory, and classical mechanics must be recovered from it as an approximation. One may use a classical Hamiltonian as a guess, but not as a derivation from something deeper.

\section{Recalling the Schr\"odinger Equation and Guessing the Free Hamiltonian}

Once the historical prologue has done its work, the lecture pivots sharply back: let us get back to the Schr\"odinger equation and really go through it. The general time-evolution law is
\begin{equation}
i\hbar \frac{\partial}{\partial t}\ket{\psi} = \hat H \ket{\psi},
\end{equation}
and the momentum operator on a line is
\begin{equation}
\hat P = -\,i\hbar \frac{\partial}{\partial x}.
\end{equation}
Susskind stresses a point that matters for the whole chapter. The relation between time derivative and Hamiltonian, and the relation between space derivative and momentum, are part of the quantum framework itself. They are not derived from classical mechanics. The minus sign in the momentum operator is a convention, but it is the convention we keep.

Now we specialize to a nonrelativistic particle moving on a line. Here the lecture openly imports a classical formula, but only as a guess:
\begin{equation}
\hat H = \frac{\hat P^2}{2m}.
\end{equation}
Why \(\hat P^2\) and not \(\hat P\)? Because momentum carries a direction, while the energy of a free nonrelativistic particle should not change when we reverse that direction. A particle moving left and a particle moving right with the same magnitude of momentum should have the same energy. That already forces the dependence to be even in \(p\), and the simplest choice is \(p^2\). The factor \(2m\) is then the familiar nonrelativistic normalization.

This is a characteristic Susskind move: first make the guess, then ask what it means. In particular, what does \(\hat P^2\) mean? It simply means that the operator \(\hat P\) acts twice.

\section{From Kets to Wave Functions: The Free-Particle Equation}

At this point the lecture shifts from abstract state-vector language to wave functions on a line. The board snapshot below records that change in layout: the board shows only an isolated ket and an arrow, but the spoken lecture is already moving toward the wave function \(\psi(x,t)\), so the figure should be read as evidence for the transition rather than as a complete formula.

\begin{figure}[h]
\centering
\includegraphics[width=0.42\textwidth]{lecture_09_figure_02.png}
\caption{Ket state with rightward arrow}
\end{figure}

The logical point is delicate and worth preserving. \(\psi\) depends on \(x\) because position is an observable. But \(\psi\) also changes with time, and that dependence is of a different character. One does not ask for the probability that the time is such-and-such; one asks for the probability that the particle is at position \(x\) at a given time \(t\). Thus \(\psi(x,t)\) is a time-evolving state written in the position representation.

Now substitute the free-particle Hamiltonian into the general Schr\"odinger equation:
\begin{align}
i\hbar \frac{\partial}{\partial t}\psi(x,t)
&= \frac{1}{2m}\hat P^2 \psi(x,t) \\
&= \frac{1}{2m}\left(-i\hbar \frac{\partial}{\partial x}\right)\left(-i\hbar \frac{\partial}{\partial x}\right)\psi(x,t).
\end{align}
Since \((-i)^2=-1\), this becomes
\begin{equation}
i\hbar \frac{\partial}{\partial t}\psi(x,t)
=
-\,\frac{\hbar^2}{2m}\frac{\partial^2}{\partial x^2}\psi(x,t).
\end{equation}
This is the free-particle Schr\"odinger equation.

\paragraph{Worked derivation.}
The lecture does not present this as a theorem dropped from the sky. It unfolds the substitution line by line. First write the time derivative on the left. Then replace \(\hat H\) by \(\hat P^2/2m\). Then remember that \(\hat P\) means \(-i\hbar \partial_x\). Acting twice gives a second derivative in \(x\). Only after that algebra does the standard equation emerge. The point is to see how the operator equation becomes a differential equation in position space.

\section{Plane Waves, Momentum Eigenstates, and the Meaning of Energy}

The next question is exactly the one the lecture asks: what are the solutions? Susskind does not begin with the most general answer. He first looks for a trial solution that oscillates in space and in time,
\begin{equation}
\psi(x,t) = e^{ipx/\hbar} e^{i\omega t}.
\end{equation}
The choice of \(e^{ipx/\hbar}\) is not arbitrary. Those are the wave functions that are eigenfunctions of the momentum operator. The factor \(e^{i\omega t}\) is then added to see what time dependence the Schr\"odinger equation allows.

Substitute this ansatz into the free-particle equation. Differentiating with respect to time pulls down \(i\omega\), so the left-hand side becomes
\begin{equation}
i\hbar \frac{\partial \psi}{\partial t} = i\hbar (i\omega)\psi = -\,\hbar\omega\,\psi.
\end{equation}
Differentiating twice with respect to \(x\) pulls down \((ip/\hbar)^2=-p^2/\hbar^2\), so the right-hand side becomes
\begin{equation}
-\,\frac{\hbar^2}{2m}\frac{\partial^2 \psi}{\partial x^2}
=
-\,\frac{\hbar^2}{2m}\left(-\,\frac{p^2}{\hbar^2}\right)\psi
=
\frac{p^2}{2m}\psi.
\end{equation}
Cancelling the common factor \(\psi\), we find
\begin{equation}
-\hbar\omega = \frac{p^2}{2m},
\qquad\text{so}\qquad
\omega = -\,\frac{p^2}{2m\hbar}.
\end{equation}
Therefore the plane-wave solution takes the form
\begin{equation}
\psi_p(x,t)
=
e^{ipx/\hbar - i p^2 t/(2m\hbar)}.
\end{equation}

This is one of the lecture's main payoffs. The \(x\)-dependence identifies the state as a momentum eigenstate with momentum \(p\), while the Hamiltonian eigenvalue is
\begin{equation}
E = \frac{p^2}{2m}.
\end{equation}
So the free particle energy is the familiar nonrelativistic expression, not because classical mechanics has been inserted as law, but because the quantum guess \(\hat H=\hat P^2/2m\) produces states whose time dependence is exactly that of energy eigenstates with that value of \(E\).

There is immediately another state with the same energy:
\begin{equation}
\psi_{-p}(x,t)
=
e^{-ipx/\hbar - i p^2 t/(2m\hbar)}.
\end{equation}
Its momentum is opposite, but the time dependence is identical because only \(p^2\) enters the energy.

\subsection{Question \& Answer}

If the energy is \(E\), why do two momentum states contribute, and how do we compute the probability for that energy?

Because the free-particle energy depends on \(p^2\), both \(p\) and \(-p\) produce the same value of \(E\). If we form a superposition
\begin{equation}
\psi(x,t)=\alpha\,\psi_p(x,t)+\beta\,\psi_{-p}(x,t),
\end{equation}
then the probability for right-moving momentum is \(\alpha^*\alpha\), while the probability for left-moving momentum is \(\beta^*\beta\). But both components carry the same energy \(E=p^2/(2m)\). Thus the lecture's rule is simple: to get the probability associated with that energy, we add the contributions from the two momentum states that share it. The conceptual point is the important one. Momentum labels and energy labels are not the same labels.

\section{General Solution in Momentum Space and the Fourier Transform at \(t=0\)}

Only now does the lecture ask for the general solution. The answer is to superpose the special plane-wave solutions over all momenta:
\begin{equation}
\psi(x,t)
=
\int_{-\infty}^{\infty}
\frac{dp}{\sqrt{2\pi}}\,
e^{i p x/\hbar - \frac{i p^2}{2m\hbar}t}\,
\tilde\psi(p).
\end{equation}
Here \(\tilde\psi(p)\) is the amplitude for momentum \(p\). Since the integral runs from \(-\infty\) to \(+\infty\), both right-moving and left-moving components are automatically included. In general, this is a superposition of different momenta and also of different energies.

The lecture then asks a practical question: what is the probability density for momentum? Since momentum is continuous, we should ask for a density, not for the probability of one isolated value. The board snapshot below captures the moment when the notation is corrected and the answer is written with a new label \(q\), precisely because \(p\) is already the dummy variable of integration.

\begin{figure}[h]
\centering
\includegraphics[width=0.55\textwidth]{lecture_09_figure_03.png}
\caption{Momentum-space probability and free-particle energy}
\end{figure}

The corresponding formulas are
\begin{equation}
P(q)=\tilde\psi^*(q)\tilde\psi(q),
\end{equation}
and
\begin{equation}
E=\frac{p^2}{2m}.
\end{equation}
The two equations belong together. The first tells us how momentum is distributed in the state; the second tells us how those momentum labels are converted into energy labels.

If we ask instead for the contribution from a fixed energy \(E\), the lecture proceeds in the simplest way. First solve
\begin{equation}
E=\frac{p^2}{2m},
\end{equation}
which gives
\begin{equation}
p_\pm = \pm \sqrt{2mE}.
\end{equation}
Then the fixed-energy contribution comes from both momentum values:
\begin{equation}
P(E)
=
\tilde\psi^*(p_+)\tilde\psi(p_+)
+
\tilde\psi^*(p_-)\tilde\psi(p_-).
\end{equation}
What matters here is not a refined density formula, but the fact that one energy generally receives contributions from two opposite momenta.

The next question is what \(\tilde\psi(p)\) actually means. Susskind answers this by stepping to the initial time \(t=0\). On the board he temporarily suppresses \(\hbar\) in the Fourier-transform formulas, and it is worth preserving that shorthand because it is exactly the board rhythm recorded by the next figure.

\begin{figure}[h]
\centering
\includegraphics[width=0.55\textwidth]{lecture_09_figure_04.png}
\caption{Momentum integral rewritten as the initial wave function}
\end{figure}

At \(t=0\), the general superposition becomes
\begin{equation}
\psi(x,0)
=
\int_{-\infty}^{\infty}
\frac{dp}{\sqrt{2\pi}}\,
e^{ipx}\,\tilde\psi(p),
\end{equation}
and in the lecture this is immediately shortened to \(\psi(x)\) when one is speaking only about the initial data. The inverse transform is
\begin{equation}
\tilde\psi(p)
=
\int_{-\infty}^{\infty}
\frac{dx}{\sqrt{2\pi}}\,
e^{-ipx}\,\psi(x,0).
\end{equation}
Thus \(\tilde\psi(p)\) is not an arbitrary coefficient function invented for convenience. It is the Fourier transform of the initial wave function.

\subsection{Question \& Answer}

Could this free-particle equation describe a photon?

No. The lecture's answer is immediate and precise. This is a nonrelativistic theory. Its Hamiltonian is built from
\begin{equation}
E=\frac{p^2}{2m},
\end{equation}
and photons do not satisfy that relation. They do not even carry a mass parameter \(m\) of this sort. So the equation is not a photon equation. It is the appropriate description for particles moving well below the speed of light.

\section{Dispersive Time Evolution from Initial Data}

Once the Fourier-transform relation is in hand, the free-particle evolution becomes an algorithm. One begins with the initial wave function \(\psi(x,0)\), computes its Fourier transform \(\tilde\psi(p)\), and then reinstates the time-dependent phase for each momentum component. In effect,
\begin{equation}
\tilde\psi(p,t)
=
e^{-i p^2 t/(2m\hbar)}\,\tilde\psi(p,0),
\end{equation}
so that
\begin{equation}
\psi(x,t)
=
\int_{-\infty}^{\infty}
\frac{dp}{\sqrt{2\pi}}\,
e^{i p x/\hbar}\,
e^{-i p^2 t/(2m\hbar)}\,
\tilde\psi(p,0).
\end{equation}

The lecture presents this not as a formal identity but as a practical update rule:
\begin{enumerate}
\item Specify the initial wave function \(\psi(x,0)\).
\item Compute its Fourier transform \(\tilde\psi(p)\).
\item Multiply each momentum component by the phase \(e^{-i p^2 t/(2m\hbar)}\).
\item Reconstruct \(\psi(x,t)\) by the inverse transform.
\item Then compute \(|\psi(x,t)|^2\) to obtain the position probability density at time \(t\).
\end{enumerate}

This is also where the lecture emphasizes that free-particle evolution is dispersive. The shape of a generic wave packet is not preserved. Different Fourier components carry different values of \(p\), hence different energies, hence different time phases. The packet therefore changes shape as it evolves. The lecture briefly acknowledges the language of phase velocity and group velocity, but does not pursue it. The essential point for this chapter is simpler: different momentum components evolve differently, so a general wave packet does not move rigidly.

\section{Expectation Values, Commutators, Potentials, and the Classical Limit}

After the free-particle story is in place, the lecture backs off and changes scale. The real goal is to understand how classical mechanics emerges, at least in the motion of averages. Susskind therefore returns to a general observable \(\hat K\) and studies the time evolution of its expectation value.

\paragraph{Expectation values.}
Start with
\begin{equation}
\langle K\rangle = \langle\psi|\hat K|\psi\rangle.
\end{equation}
The operator \(\hat K\) is taken to be time independent; only the bra and ket change with time. Differentiating and using the Schr\"odinger equation gives
\begin{align}
\frac{d}{dt}\langle K\rangle
&=
\langle \dot\psi|\hat K|\psi\rangle
+
\langle \psi|\hat K|\dot\psi\rangle \\
&=
\frac{i}{\hbar}\langle \psi|\hat H \hat K|\psi\rangle
-
\frac{i}{\hbar}\langle \psi|\hat K \hat H|\psi\rangle \\
&=
\frac{i}{\hbar}\langle [\hat H,\hat K]\rangle.
\end{align}
On the board the lecture temporarily sets \(\hbar=1\) to keep the algebra moving, but the formula above is the version with \(\hbar\) restored.

The classical analog is
\begin{equation}
\dot K = \{K,H\}.
\end{equation}
The lecture pauses over signs and ordering here, but the point is that classical time evolution is controlled by the Poisson bracket, while quantum time evolution of expectation values is controlled by the commutator.

\paragraph{Poisson brackets and commutators.}
The simplest example is a function of position and the momentum. Classically,
\begin{equation}
\{f(x),p\}=f'(x),
\qquad
\{x,p\}=1.
\end{equation}
Quantum mechanically one computes the commutator by letting it act on an arbitrary wave function:
\begin{align}
[f(\hat X),\hat P]\psi(x)
&=
f(x)\left(-i\hbar\frac{d\psi}{dx}\right)
-
\left(-i\hbar\frac{d}{dx}\right)\bigl(f(x)\psi(x)\bigr) \\
&=
f(x)\left(-i\hbar\frac{d\psi}{dx}\right)
+
i\hbar\left(f(x)\frac{d\psi}{dx}+f'(x)\psi(x)\right) \\
&=
i\hbar f'(x)\psi(x).
\end{align}
Hence
\begin{equation}
[f(\hat X),\hat P]=i\hbar f'(\hat X).
\end{equation}
By the reciprocal role of \(x\) and \(p\), one likewise has
\begin{equation}
[\hat X,f(\hat P)]=i\hbar f'(\hat P),
\end{equation}
and in particular
\begin{equation}
[\hat X,\hat P]=i\hbar.
\end{equation}

The lecture then checks that Poisson brackets and commutators share the same algebraic structure. Both are antisymmetric:
\begin{equation}
\{A,B\}=-\{B,A\},
\qquad
[A,B]=-[B,A].
\end{equation}
Both satisfy a product rule. The corrected quantum identity, after the board correction, is
\begin{align}
\{AB,C\} &= A\{B,C\}+B\{A,C\}, \\
[AB,C] &= A[B,C]+[A,C]B.
\end{align}
Thus, in the simple situations relevant here, the lecture's bridge is
\begin{equation}
[A,B]\sim i\hbar\,\{A,B\}.
\end{equation}
The two structures are not always identical in complicated operator-ordering problems, but they are close enough to reveal how classical mechanics emerges.

\paragraph{Particle in a potential.}
Now we return to a particle on a line, but with a potential:
\begin{equation}
\hat H = \frac{\hat P^2}{2m}+U(\hat X).
\end{equation}
In position space this gives the Schr\"odinger equation
\begin{equation}
i\hbar\frac{\partial \psi(x,t)}{\partial t}
=
\left[
-\,\frac{\hbar^2}{2m}\frac{\partial^2}{\partial x^2}
+
U(x)
\right]\psi(x,t).
\end{equation}
The new term \(U(x)\psi(x,t)\) is not an extra number added to the equation. It is the action of the multiplication operator \(U(\hat X)\) on the state.

We now apply the expectation-value formula first to position. Since any function of \(\hat X\) commutes with \(\hat X\),
\begin{equation}
[U(\hat X),\hat X]=0.
\end{equation}
Thus only the kinetic term contributes:
\begin{equation}
\frac{d}{dt}\langle X\rangle
=
\frac{i}{\hbar}
\left\langle
\left[
\frac{\hat P^2}{2m},\hat X
\right]
\right\rangle.
\end{equation}
Using the corrected product rule,
\begin{equation}
[\hat P^2,\hat X]
=
\hat P[\hat P,\hat X]+[\hat P,\hat X]\hat P
=
-2i\hbar \hat P.
\end{equation}
Therefore
\begin{equation}
\frac{d}{dt}\langle X\rangle
=
\frac{\langle P\rangle}{m}.
\end{equation}
So the average velocity is the average momentum divided by the mass, exactly as one would hope.

Now apply the same rule to momentum:
\begin{equation}
\frac{d}{dt}\langle P\rangle
=
\frac{i}{\hbar}
\left\langle
\left[
\frac{\hat P^2}{2m}+U(\hat X),\hat P
\right]
\right\rangle.
\end{equation}
Here \(\hat P\) commutes with \(\hat P^2\), so only the potential matters:
\begin{equation}
\frac{d}{dt}\langle P\rangle
=
\frac{i}{\hbar}\langle [U(\hat X),\hat P]\rangle.
\end{equation}
Using \([f(\hat X),\hat P]=i\hbar f'(\hat X)\), we obtain
\begin{equation}
\frac{d}{dt}\langle P\rangle
=
-\,\langle U'(\hat X)\rangle.
\end{equation}
If we define the force by
\begin{equation}
F(x)=-U'(x),
\end{equation}
then
\begin{equation}
\frac{d}{dt}\langle P\rangle = \langle F(\hat X)\rangle.
\end{equation}

These are the exact quantum equations for the expectation values:
\begin{equation}
\frac{d}{dt}\langle X\rangle=\frac{\langle P\rangle}{m},
\qquad
\frac{d}{dt}\langle P\rangle=-\,\langle U'(\hat X)\rangle.
\end{equation}
Combining them gives
\begin{equation}
m\frac{d^2}{dt^2}\langle X\rangle = -\,\langle U'(\hat X)\rangle.
\end{equation}

Now comes the controlled classical limit. If the wave packet is well localized, and if the potential is smooth on the scale over which the wave function oscillates, then the force does not vary much across the packet. In that case
\begin{equation}
\langle F(\hat X)\rangle \approx F(\langle X\rangle),
\end{equation}
and therefore
\begin{equation}
m\frac{d^2}{dt^2}\langle X\rangle \approx F(\langle X\rangle).
\end{equation}
This is the sense in which the center of the packet follows classical mechanics. The exact quantum statement involves \(\langle U'(\hat X)\rangle\), not \(U'(\langle X\rangle)\). The latter is an approximation that becomes good for smooth potentials and coherent, localized packets.

\section{Summary}

The lecture begins by insisting that quantum mechanics was forced on physics not by taste but by crisis: blackbody radiation, photons, atomic spectra, and finally probabilistic decay made the classical worldview untenable. Only after that conceptual clearing does the mathematics resume.

For a nonrelativistic particle on a line, the guess \(\hat H=\hat P^2/2m\) leads to the free-particle Schr\"odinger equation, plane-wave momentum eigenstates, and the momentum-space superposition formula. The Fourier transform then turns initial data into later-time evolution, and that evolution is generically dispersive. Finally, by stepping back to expectation values and commutators, the lecture recovers the exact equations
\[
\frac{d}{dt}\langle X\rangle=\frac{\langle P\rangle}{m},
\qquad
\frac{d}{dt}\langle P\rangle=-\,\langle U'(\hat X)\rangle,
\]
which show how classical motion emerges from quantum mechanics in the appropriate limit. Classical mechanics is not the starting point. It is the late approximation.